% Conclusion chapter for:
% A Hierarchical Time-Aware Transformer for Flow-Level Network Intrusion Detection
%
% The chapter number is assigned automatically by LaTeX. When this file follows
% Chapter 8 in the thesis, it will be rendered as Chapter 9.

\chapter{Conclusion and Future Work}
\label{sec:conclusion_future_work}

This thesis investigated whether network intrusion detection can benefit from
preserving two forms of structure that are commonly weakened by flow-level
aggregation: the temporal organization of packets within a flow and the causal
relationships among flows. To address this question, the study developed and
evaluated a hierarchical, payload-independent detection framework. Stage~1
learns an intra-flow representation from ordered packet metadata, explicit
elapsed-time information, and aggregate flow statistics. Stage~2 then refines
the decision by attending to earlier flows selected through a specified network
relation. The resulting design separates representation learning from
contextual reasoning, allowing the contribution of each stage to be evaluated
under controlled comparisons.

The experiments provide consistent evidence for the central thesis: preserving
temporal and hierarchical traffic structure improves minority-class detection
without requiring direct payload bytes. This conclusion is supported not only
by accuracy or ROC-AUC, which can be optimistic under class imbalance, but also
by macro-F1, macro-precision, macro-recall, attack-class F1, PR-AUC, false
positives, false negatives, and FPR. At the same time, the conclusions remain
bounded by the organizational capture, binary label construction, and
single-seed evidence described in Chapters~\ref{sec:dataset_preprocessing},
\ref{sec:experimental_methodology}, and~\ref{sec:results}.

\section{Answers to the Research Questions}
\label{subsec:conclusion_rq_answers}

\paragraph{RQ1: Does time-aware packet-sequence modelling produce more
informative representations than flow-level aggregation?}
The controlled Stage~1 experiments support an affirmative answer. On the
35,410-flow test manifest, the Transformer using both packet position and
cumulative time achieved macro-F1 $0.8882$ and PR-AUC $0.8564$, compared with
$0.7762$ and $0.5604$, respectively, for the flow-statistics MLP. The sequence
model also reduced FP from 853 to 296 and FN from 501 to 306, lowering FPR from
$2.5086\%$ to $0.8705\%$. Removing either positional or temporal encoding
reduced macro-F1, macro-recall, and attack-class F1, showing that packet order
and elapsed time provide complementary rather than interchangeable evidence.

The fusion study further showed that packet and flow views are most useful when
they interact in a controlled manner. Scheme~C, which applies bounded
flow-conditioned token modulation before Transformer encoding and retains late
fusion, reached macro-F1 $0.8955$ and PR-AUC $0.8763$ on the same controlled
manifest. The sequence-length and selection studies qualified this result:
$L=128$ with deterministic head selection produced the strongest overall error
balance, whereas increasing the packet budget to 256 did not improve
performance. Thus, the answer to RQ1 is not simply that longer sequences are
better; rather, an appropriate packet horizon, explicit time, packet order, and
flow-level conditioning jointly produce the most informative representation.

\paragraph{RQ2: Does causal cross-flow context improve detection beyond
independent per-flow classification?}
The Stage~2 experiments also support an affirmative answer. On the common
84,380-flow test manifest, the independent Scheme~C classifier achieved
macro-F1 $0.8954$, PR-AUC $0.8763$, and FPR $0.8249\%$. Adding causal
source-host context increased macro-F1 to $0.9275$ and PR-AUC to $0.9345$ while
reducing FPR to $0.5429\%$. In absolute terms, FP decreased from 667 to 439 and
FN from 732 to 527. Because both error types declined on the same test flows,
the improvement cannot be explained only as a threshold-dependent exchange
between precision and recall.

Source-host context with a 128-flow window provided the best balanced operating
point among the tested relation and window definitions. Endpoint context
offered higher attack recall at the cost of more false alarms, indicating that
the preferred relation may depend on operational priorities. The proposed
target-query gated attention model also exceeded the no-context, vanilla
Transformer, LSTM, GRU, and CNN--LSTM controls in macro-F1 and PR-AUC. These
findings show that useful intrusion evidence can be distributed across related
flows and that relation-specific, causal context is more effective than treating
each flow independently.

\section{Principal Contributions}
\label{subsec:conclusion_contributions}

The thesis makes the following contributions.

\begin{enumerate}
    \item It establishes an end-to-end, payload-independent data path from PCAP
    capture through a custom Suricata output module to aligned packet- and
    flow-level records. The construction preserves packet order, direction,
    timing, flow identity, and aggregate statistics required by the two-stage
    model.

    \item It introduces a Stage~1 representation that combines explicit
    cumulative-time encoding, conventional packet-position encoding,
    mask-aware Transformer processing, hybrid pooling, and bounded
    flow-conditioned token modulation. The accompanying ablations distinguish
    the effects of temporal information, packet order, statistical fusion,
    sequence length, and packet selection.

    \item It introduces a causal Stage~2 mechanism in which the target flow
    queries a relation-specific history through gated multi-head attention.
    This formulation supports time-only, source-host, destination-host, and
    endpoint contexts while preventing future-flow information from entering
    the current prediction.

    \item It provides an imbalance-aware evaluation that treats macro-F1 and
    PR-AUC as primary outcomes and interprets them jointly with class-1
    performance, confusion matrices, and FPR. Test-manifest boundaries and
    validation-selected thresholds are kept explicit so that only controlled
    results are compared directly.
\end{enumerate}

Together, these contributions are more than an application of a Transformer to
tabular network data. They define and test a hierarchical representation of
traffic in which packet dynamics explain individual flows and prior related
flows provide decision context.

\section{Limitations and Scope of the Conclusions}
\label{subsec:conclusion_limitations}

The validity boundaries have been discussed alongside the experimental results;
only those that materially constrain the final claims are summarized here. The
dataset originates from one organizational PCAP capture and uses alert-derived
binary labels. Consequently, the experiments demonstrate effectiveness on the
observed environment but do not yet establish generalization across networks,
time periods, capture tools, attack families, or labeling policies. The model
uses no direct payload bytes, but this alone does not prove robustness across
all encrypted protocols or previously unseen attacks.

Most neural experiments were conducted with one random seed. Large performance
gaps provide useful empirical evidence, whereas small differences between
nearby configurations should remain descriptive until repeated runs and paired
uncertainty estimates are available. In addition, Stage~1 and Stage~2 were
trained sequentially using frozen exported embeddings; this improves
experimental separation but may prevent end-to-end optimization of the final
detection objective. Finally, the reported Stage~2 latency covers model
evaluation after embeddings and context indices are available. It is therefore
not an end-to-end measurement from PCAP ingestion to alert generation, and the
current study does not quantify production memory use, tail latency, or concept
drift.

\section{Future Work}
\label{subsec:conclusion_future_work}

The empirical findings motivate two principal methodological extensions:
resource-efficient pre-training for Stage~1 and attack-adaptive,
relation-aware context modelling for Stage~2. These directions should be
developed together with stronger validation rather than being assessed only by
additional single-run improvements.

\subsection{Resource-Efficient Self-Supervised Pre-Training}
\label{subsubsec:future_pretraining}

The current Stage~1 encoder is trained from supervised binary labels. This is
practical for the available experiment, but it limits representation learning
to the information exposed by a relatively small set of alert-derived labels.
Resource constraints prevented large-scale pre-training in the present study.
A primary direction for future work is therefore to pre-train the Stage~1
encoder on a larger collection of unlabeled packet and flow records before
fine-tuning it for intrusion detection. Self-supervised traffic pre-training has
the potential to learn reusable structural representations without requiring
payload inspection or a label for every flow
\cite{2024MAETraffic,Dong2025PretrainingEncrypted}.

Several pre-training objectives are compatible with the current feature space.
Masked packet-feature modelling could reconstruct selected header, direction,
length, and timing attributes. Predictive objectives could estimate the next
packet's direction, discretized length, or log inter-arrival time. Temporal-order
verification could distinguish correctly ordered packet sequences from
perturbed sequences, while contrastive learning could align two valid
subsequences or augmented views of the same flow. Identity-like fields such as
raw IP addresses and ports should be removed or carefully randomized during
these tasks so that the encoder learns traffic behaviour rather than memorizing
hosts.

Pre-training does not eliminate computational cost; it moves additional cost to
an offline representation-learning phase. For this reason, future experiments
should investigate mixed-precision training, gradient accumulation, moderate
sequence lengths, parameter sharing, and progressive training on increasingly
large traffic samples. Downstream evaluation should compare random
initialization, frozen-encoder linear probing, partial fine-tuning, and full
fine-tuning under equal labelled-data budgets. The key question is not only
whether pre-training raises in-domain macro-F1, but whether it improves label
efficiency, rare-attack detection, convergence speed, and transfer to later or
external captures.

\subsection{Attack-Adaptive Relation-Aware Context Modelling}
\label{subsubsec:future_relation_aware_context}

The Stage~2 experiments show that context quality depends on how related flows
are defined. Source-host context provides the best aggregate error balance,
whereas endpoint context produces higher attack recall with more false
positives. This suggests that no single relation is necessarily optimal for all
attack mechanisms. Scanning may be represented by one source contacting many
destinations, distributed denial-of-service activity by many sources contacting
one destination, brute-force behaviour by repeated service interactions, and
lateral movement by a chain of temporally connected endpoints. Future Stage~2
models should therefore select context according to the target flow and the
behaviour being modelled rather than relying on one globally fixed relation.

A relation-aware Transformer could construct one causal context pool containing
source-host, destination-host, endpoint, protocol, service, subnet, and
time-neighbour relations. Each context item would receive a learned relation
embedding or relation-specific attention bias. Alternatively, separate
key--value projections could be learned for each relation and combined through
a target-conditioned gate. In either design, only flows that finish before the
target begins should remain eligible, preserving the causal constraint of the
current system. The context horizon could also be adaptive: a model could choose
between a time-duration window and a flow-count window according to recent
traffic density.

Evaluation of this extension requires attack-family labels or a defensible
behavioural grouping in addition to the current binary target. Results should
report per-family relation weights, recall, FPR, and confusion patterns to test
whether different attacks actually use different contexts. Comparisons must
include the current single-relation model, a relation-agnostic Transformer, and
parameter-matched multi-relation controls. Host identifiers should be masked in
cross-time and cross-site tests so that relation-aware attention cannot obtain
its apparent gain by memorizing machines associated with historical alerts.
Relation weights may help analysis, but they should not be treated as causal
explanations without relation ablation and counterfactual context removal.

\subsection{Statistical and External Validation}
\label{subsubsec:future_validation}

The principal Stage~1 and Stage~2 comparisons should be repeated with at least
three to five random seeds. Each run should archive the configuration,
validation-selected threshold, checkpoint, test labels, probabilities, and
predictions. These aligned outputs would support confidence intervals, paired
bootstrap comparisons, McNemar tests, calibration analysis, and genuine
precision--recall curves. The models should then be evaluated on
chronologically later captures, unseen hosts, public benchmarks, and, where
governance permits, other organizational networks. This would distinguish
in-domain performance from temporal, cross-site, and attack-family transfer.

\subsection{Streaming Deployment, Efficiency, and Robustness}
\label{subsubsec:future_streaming_deployment}

The causal formulation is compatible with online inference, but a streaming
implementation remains to be validated. Future work should maintain bounded
per-relation buffers, evict expired flows, handle out-of-order records, and
define behaviour under packet loss, timestamp noise, feature missingness, NAT,
and highly shared servers. Evaluation should cover the full path from PCAP
processing and Suricata extraction through both model stages and alert
generation. In addition to predictive performance, it should report throughput,
peak memory, and median, 95th-percentile, and 99th-percentile latency under
realistic load.

\section{Closing Remarks}
\label{subsec:conclusion_closing}

The results of this thesis show that network traffic should not be treated only
as independent rows of aggregate statistics. Packet-level timing and ordering
carry information that is lost during aggregation, and related flows provide
additional evidence that cannot be recovered from an isolated flow. A
hierarchical model that preserves both structures achieved a better balance
between detecting the positive class and controlling false alarms in the
studied environment. The next step is therefore not merely to enlarge the
model, but to verify this structural advantage across seeds, captures, and
operational conditions. Resource-efficient pre-training can then strengthen the
flow representation, while attack-adaptive relation-aware attention can extend
context modelling beyond one fixed relation. Together, these directions provide
a practical research path towards more transferable, context-aware, and
payload-independent network intrusion detection.
